\subsection{代数的基。乘法表}

根据定义，A-代数 E 的基就是 E 的 A-模结构的基。设 \((a_{i})_{i\in I}\) 是 E 的一组基；存在唯一的由环 A 的元素组成的族 \((\gamma_{ij}^{k})_{(i,j,k)\in I\times I\times I}\) ，使得对每个有序对 \((i,j)\in I\times I\) ， \(k\in I\) 中使 \(\gamma_{ij}^{k}\neq 0\) 者只有有限多个，且

\[
a _ {i} a _ {j} = \sum_ {k \in \mathbf {I}} \gamma_ {i j} ^ {k} a _ {k}.\tag{7}
\]

诸 \(\gamma_{ij}^{k}\) 称为代数 E 关于基 \((a_{i})\) 的结构常数，而关系式 (7) 构成代数 E （相对于基 \((a_{i})\) ）的乘法表。

关系式（7）可以想象为将这些关系式的右端排列在一个方形表格中写出。

\(\cdots\)

\(a_j\)

\(\cdots\)

\(\vdots\)

\(a_j\)

\(\sum_{k} \gamma_{ij}^{k} a_k\)

\(\vdots\)

这里约定，第 i 行第 j 列处的元素等于乘积 \(a_{i}a_{j}\) 。

反之，给定一个 A-模 E、E 的基 \((a_{i})_{i\in I}\) 以及 A 的元素的一个族 \((\gamma_{ij}^{k})\) ，使得对每个有序对 \((i,j)\in I\times I\) ， \(k\in I\) 中使 \(\gamma_{ij}^{k}\neq 0\) 者只有有限多个，则在 E 上存在唯一的 A-代数结构使关系式 (7) 成立，因为 A-模 \(\mathbf{E}\otimes_{\mathbf{A}}\mathbf{E}\) 是自由的，且以 \((a_{i}\otimes a_{j})_{(i,j)\in I\times I}\) 为基（参见 II，§ 3，第 7 号，命题 7 的推论 2）。

设 E 是一个 A-代数， \((a_{i})_{i\in I}\) 是 A-模 E 的一个生成系（例如一组基）。要使 E 为结合的，必要且充分的条件是诸 \(a_{i}\) 满足结合性关系

\[
(a _ {i} a _ {j}) a _ {k} = a _ {i} (a _ {j} a _ {k}) \quad \text { for   all } i, j, k\tag{8}
\]

映射 \((x, y, z) \mapsto (xy)z - x(yz)\) 是一个 A-三线性映射

\[
\mathbf {E} \times \mathbf {E} \times \mathbf {E} \rightarrow \mathbf {E}
\]

从而定义了一个 A-线性映射 \(E \otimes_{A} E \otimes_{A} E \to E\) ；若后一映射对所有元素 \(a_{i} \otimes a_{j} \otimes a_{k}\) （它们构成 A-模 \(E \otimes_{A} E \otimes_{A} E\) 的一个生成系）均为零，则它恒为零。

类似地，要使 E 为交换的，必要且充分的条件是诸 \(a_{i}\) 满足交换性关系

\[
a _ {i} a _ {j} = a _ {j} a _ {i} \quad \text { for   all } i, j.\tag{9}
\]

证明是类似的，这次考虑 A-双线性映射 \((x,y)\mapsto xy - yx\) 。最后，要使元素 \(e\in E\) 成为单位元，必要且充分的条件是诸 \(a_{i}\) 满足关系

\[
a _ {i} = e a _ {i} = a _ {i} e \quad \text { for   all } i,\tag{10}
\]

这一次是借助 A-线性映射 \(x \mapsto x - xe\) 与 \(x \mapsto x - ex\) 看出的。

当 \((a_{i})_{i\in I}\) 是 E 的一组基且 \((\gamma_{ij}^{k})\) 是相应的结构常数族时，关系式 (8) 等价于关系式

\[
\sum_ {r} \gamma_ {i j} ^ {r} \gamma_ {r k} ^ {s} = \sum_ {r} \gamma_ {i r} ^ {s} \gamma_ {j k} ^ {r}
\]

对所有 \(i, j, k, s\) 成立。类似地，关系式 (9) 等价于 \(\gamma_{ij}^{k} = \gamma_{ji}^{k}\) （对所有 \(i, j, k\) ）。

设 \((a_{i})_{i\in I}\) 是 A-代数 E 的一组基；若 \(\rho :A\to B\) 是一个环同态，则 \((a_{i}\otimes 1)\) 是 B-代数 \(\rho^{*}(\mathbf{E}) = \mathbf{E}\otimes_{\mathbf{A}}\mathbf{B}\) 的一组基（II，§ 5，第 1 号，命题 4）。若 \((\gamma_{ij}^{k})\) 是关于基 \((a_{i})\) 的结构常数族，则族 \((\rho (\gamma_{ij}^{k}))\) 是 \(\rho^{*}(\mathbf{E})\) 关于基 \((a_{i}\otimes 1)\) 的结构常数族。

